
Selecting the LoRA rank for each transformer layer currently requires training, gradient signals, or calibration data, yet the pretrained weights themselves encode a structural signal about how much rank each layer needs. This paper investigates whether effective rank $\text{erank}(W_0) = \exp(H(\sigma/\|\sigma\|_1))$ --- a scalar derived from the singular value entropy of each pretrained weight matrix --- predicts per-layer optimal LoRA rank without any task data. For BERT-base-uncased, erank correlates with per-layer PARA oracle ranks at Pearson $r = 0.984$ ($p = 0.0013$, one-tailed, $n = 5$ oracle-measured layers): FFN intermediate layers (erank $\approx 704--710$) receive oracle rank $r^* = 64$, while attention layers --- including output projections and query projection (erank $\approx 557--590$) --- receive oracle rank $r^* = 4$. The near-perfect correlation reflects binary discrimination between two structural tiers rather than a smooth graded relationship across the full erank range; 5 of 72 target layers have been oracle-measured, and full coverage is pending. Bootstrap confidence interval computation returned NaN values due to the degenerate bimodal sample structure; the reported $p$-value is from the exact one-tailed $t$-distribution. The participation ratio agrees with erank in layer ranking at Pearson $\rho = 0.968$. erank maps computed for BERT-base-uncased, DeBERTa-v3-base, and ViT-base-patch16-224 show consistent layer-type structure across architectures, with ViT exhibiting the largest within-model erank variation (range ratio $2.97\times$. DeBERTa-v3-base and ViT oracle sweeps are pending, so the pre-registered multi-family gate ($\geq 2/3$ families) has not been met; results are presented as preliminary single-family evidence.

